% Generated draft: requires mathematical and source review.
\documentclass[11pt]{article}
\usepackage[margin=25mm]{geometry}
\usepackage{fontspec}
\setmainfont{Latin Modern Roman}
\usepackage{amsmath,amssymb,mathtools,unicode-math}
\setmathfont{Latin Modern Math}
\usepackage{enumitem,needspace,xurl,hyperref}
\setlength{\parindent}{0pt}
\setlength{\parskip}{5pt}
\setlength{\emergencystretch}{3em}
\newcommand{\N}{\mathbb{N}}
\newcommand{\Z}{\mathbb{Z}}
\newcommand{\Q}{\mathbb{Q}}
\newcommand{\R}{\mathbb{R}}
\newcommand{\C}{\mathbb{C}}
\newcommand{\F}{\mathbb{F}}
\newcommand{\A}{\mathbb{A}}
\newcommand{\PP}{\mathbb{P}}
\newcommand{\E}{\mathbb{E}}
\newcommand{\eps}{\varepsilon}
\newcommand{\dd}{\,\mathrm{d}}
\begin{document}
\section*{UnsolvedMath ID: 30003183}
\textbf{Source problem code:} catalogue entry
\subsection{Definitions and mathematical statement}
Let \(n\in\Z\) with \(n>2\), let \(p\) be a prime, and let \(\F_p=\Z/p\Z\). Write \(D_n\) for the little \(n\)-disks operad: in positive arity \(r\), its space \(D_n(r)\) consists of configurations of \(r\) ordered, pairwise disjoint little \(n\)-disks inside the unit \(n\)-disk, with operadic composition given by insertion. A nonsymmetric cooperad retains the cocomposition maps and their compatibility identities while forgetting symmetric-group actions.

The source denotes the cochain cooperad by \(C^*(D_n;\F_p)\). For a differential graded nonsymmetric cooperad \(A\) over \(\F_p\), formality means that there is a zigzag of morphisms of differential graded nonsymmetric cooperads
\[
A\;\xleftarrow{\simeq}\;A_1\;\xrightarrow{\simeq}\;\cdots\;\xleftarrow{\simeq}\;A_m\;\xrightarrow{\simeq}\;H^*(A),
\]
where every arrow is a quasi-isomorphism in every arity and \(H^*(A)\) has zero differential and the induced cocomposition maps. Arrow directions may vary throughout the zigzag.

\textbf{Intended question.} For every integer \(n>2\) and every prime \(p\), is the full nonsymmetric cochain cooperad \(C^*(D_n;\F_p)\) formal? In the intended cochain-cooperad model, this asks for an aritywise quasi-isomorphism zigzag
\[
C^*(D_n;\F_p)\;\simeq\;H^*(D_n;\F_p)
\]
in the category of nonsymmetric differential graded cooperads, with zero differential on the right. The question concerns all arities, without an arity truncation.

\textbf{Unresolved qualifications.} The supplied record does not specify a cochain model or the construction of its cooperad structure. Contravariance of cochains gives maps into cochains on products of operad spaces; interpreting these as strict cocomposition maps into tensor products requires additional model choices or comparison data. Thus the displayed formulation is conditional on fixing the intended cochain-cooperad model. The record also does not specify whether morphisms must preserve additional multiplicative or homotopy-coherent structure, or whether nullary operations, counits, or coaugmentations are included. These conventions must be resolved before the target is fully precise; formality of the underlying nonsymmetric differential graded cooperad is the interpretation displayed above.
\subsection{Short English statement}
For every integer \(n>2\) and every prime \(p\), can the mod-\(p\) cochains of the little \(n\)-disks operad, after forgetting symmetric-group actions, be connected to their cohomology by quasi-isomorphisms preserving the intended cooperad structure in all arities? The cochain model and required structure are unspecified. The supplied dated curation labels the problem partially solved; this draft does not independently verify that assessment.
\subsection{Sources}
\begin{itemize}
\item Local UnsolvedMath registry, \texttt{data/unsolvedmath/problems.json}, permanent ID \(30003183\), source problem code \texttt{catalogue entry}, title ``Mod-p Formality of Little-Disks Cochains,'' in the Oberwolfach Reports Open Problems collection. The supplied record gives proposed year \(2016\) and names no individual proposer.
\item Supplied report citation: \emph{Topologie (2017)}, Oberwolfach Reports, DOI \texttt{10.4171/owr/2016/35}, pp.\ 2009--2068. \url{https://doi.org/10.4171/owr/2016/35}. This is also listed among the supplied literature sources as ``Oberwolfach Report 35.'' The supplied background describes these reports as records of mathematical workshops and their open-problem discussions.
\item Supplied literature source: Boavida de Brito--Horel, \url{https://arxiv.org/abs/1903.09191}. No paper title, theorem number, or precise range of applicability is provided in the record.
\item Literature assessment dated exactly \texttt{2026-08-21}: the supplied curation labels the problem \emph{partially solved} and states that higher little-disks formality in positive characteristic is known in finite ranges or truncated arities, while the full nonsymmetric statement for \(n>2\) remains unresolved. This is supplied dated curation derived from the OWR material, not an independently verified literature assessment. The record itself requests independent verification of time-sensitive status and cited claims.
\item The supplied statement status is \texttt{corrected\_verified}, with an accompanying assertion that mathematical content was preserved while notation, typography, or omitted definitions were repaired. This is source metadata, not independent mathematical verification by this draft.
\end{itemize}
\end{document}
